%% meta.tex -- Shared metadata and templates for lecture notes
%% Author: Dr.-Ing. Mark Schutera

\RequirePackage{qrcode}

%%
% End-of-document bibliography is built (so \cite{...} keys still resolve to
% author/year markers wherever they appear, including margin notes) but the
% visible thebibliography listing is suppressed: \bibitem entries are
% processed into the .aux/.bbl pipeline so cites work, then the rendered
% box is dropped. Literature thus lives only in margin notes, while the
% reference machinery stays intact.
%%
\AtBeginDocument{%
  \@ifundefined{thebibliography}{}{%
    \let\hidden@oldthebibliography\thebibliography
    \let\hidden@oldendthebibliography\endthebibliography
    \renewenvironment{thebibliography}[1]{%
      \setbox\z@\vbox\bgroup
        \hsize=\textwidth
        \hidden@oldthebibliography{##1}%
    }{%
        \hidden@oldendthebibliography
      \egroup
    }%
  }%
}

\InputIfFileExists{version.tex}{}{}
\providecommand{\docversionmajor}{0}
\providecommand{\docversionminor}{0}
\providecommand{\docversionpatch}{0}
\newcommand{\docversion}{\docversionmajor.\docversionminor.\docversionpatch}

%%
% Default author information
%%
\newcommand{\defaultauthor}{Prof. Dr.-Ing. Mark Schutera}
\newcommand{\defaultpublisher}{Unfinished Lecture Notes}

%%
% Contribution invitation. Each lecture-notes main file should set
% \notesdir to its own subdirectory, e.g.
%   \renewcommand{\notesdir}{notes_philosophyofai}
% so the QR code and URL on the copyright page point to the right book.
%%
\providecommand{\repobaseurl}{https://github.com/Quillstacks/LectureMaterial/tree/main/lecturenotes}
\providecommand{\notesdir}{}
\newcommand{\notesrepourl}{\repobaseurl/\notesdir}

\newcommand{\makecontributioninvite}{%
  \par These notes are, by their very nature, unfinished, and they improve with every reader.
  If you spot an error, disagree with a framing, or want to add a source, a question, or a position, open a pull request at \url{\notesrepourl}. Every contribution is welcome.
  \par\noindent\hfill\qrcode[height=1.6cm]{\notesrepourl}\hspace*{0pt}
}

%%
% Per-chapter version stamp. Populated by tools/update_chapter_versions.sh,
% which writes a chapter_versions.tex file with calls like
%   \setchapterversion{05_Maschinenquaelerei}{2026-05-06}{crunchy banana dachs}
% Each chapter calls \printchapterversion{<key>} right after \chapter{...}
% to emit a small italic line under the chapter title; if the key is not
% registered, the call is a silent no-op so chapter files stay safe to
% input even when the version table has not been generated yet.
%%
\newcommand{\setchapterversion}[3]{%
  \expandafter\gdef\csname chapver@date@#1\endcsname{#2}%
  \expandafter\gdef\csname chapver@name@#1\endcsname{#3}%
}
\newcommand{\printchapterversion}[1]{%
  \ifcsname chapver@date@#1\endcsname
    \par\noindent{\itshape\small\csname chapver@date@#1\endcsname{} \textperiodcentered\ \csname chapver@name@#1\endcsname}\par
    \vspace{0.25\baselineskip}%
  \fi
}
% Book-level stamp, set by tools/update_chapter_versions.sh from the
% most recent commit touching the book directory. Rendered on the
% copyright page in place of the old "Last updated, ..." line.
\newcommand{\bookver@date}{}
\newcommand{\bookver@name}{}
\newcommand{\setbookversion}[2]{%
  \renewcommand{\bookver@date}{#1}%
  \renewcommand{\bookver@name}{#2}%
}
\newcommand{\printbookversion}{%
  \ifx\bookver@date\@empty\else
    \textit{\bookver@date{} \textperiodcentered\ \bookver@name}%
  \fi
}
\InputIfFileExists{chapter_versions.tex}{}{}

%%
% Copyright page template
%%
\newcommand{\makecopyright}[1][]{%
  \newpage
  \begin{fullwidth}
  ~\vfill
  \thispagestyle{empty}
  \setlength{\parindent}{0pt}
  \setlength{\parskip}{\baselineskip}
  Copyright \copyright\ \the\year\ \thanklessauthor
  
  \par\smallcaps{Published by \thanklesspublisher}
    
\par Combobulated with the help of multiple large language model driven tools. Licensed under the Creative Commons Attribution-NonCommercial 4.0 International License (``CC BY-NC-SA 4.0''). 
You may not use this file for commercial purposes. If you remix, transform, or build upon the material, you must distribute your contributions under the same license as the original
You must obtain explicit permission from the author for uses beyond those permitted by this license. 
To view a copy of this license, visit \url{https://creativecommons.org/licenses/by-nc-sa/4.0/}. 
Unless required by applicable law or agreed to in writing, distributed material is provided on an 
\smallcaps{``AS IS'' BASIS, WITHOUT WARRANTIES OR CONDITIONS OF ANY KIND}, either express or implied. 
See the license for details.\index{license}

  \makecontributioninvite

  \par\printbookversion
  \end{fullwidth}
}

%%
% Standard front matter structure
%%
\newcommand{\makestandardfrontmatter}{%
  \frontmatter
  \maketitle
  \makecopyright
}

%%
% Standard dedication page template
%%
\newcommand{\makededication}[1]{%
  \cleardoublepage
  ~\vfill
  \begin{doublespace}
  \noindent\fontsize{18}{22}\selectfont\itshape
  \nohyphenation
  #1
  \end{doublespace}
  \vfill
  \vfill
}

%%
% Standard introduction template
%%
\newcommand{\makeintroduction}[1]{%
  \cleardoublepage
  \chapter*{Introduction}
  #1
}

%%
% Philosophy of education
%%
\newcommand{\makephilosophy}{%
  \newpage
  \chapter*{A Philosophy of Education}
  \newthought{The beautiful risk of education} is that it is an encounter.
  I am trying to model a way of thinking, and reasoning, then stepping back.
  The rest is yours.

  \vspace{2.5em}



  \begin{itemize}
    \item \textbf{Start with why.} A method understood only procedurally is a method you will misapply when it matters.
    Before learning how something works, understand why it was built, what it replaced, and where it fails.
    The mechanics then follow from necessity, not memorization.
    \item \textbf{Build taste alongside competence.} You need a sense of what beautiful is,
    you need to train your aesthetics before you can reliably produce it.
    This develops by studying what others have built, making it your own, then making your own, and honestly evaluating the gap between them.
    Taste is the compass. Skill is the engine. You need both, in that order.
    \item \textbf{Failure is the designed medium.} Each attempt teaches what a clean run cannot.
    The insight compounds across tries, not from a single correct execution.
    Those who never start failing will fail.
    \item \textbf{Work in the open.} Showing half-finished thinking, naming mistakes as they happen, is a professional skill.
    We practice it here. It is how knowledge actually moves between people. And it is how knowledge moves towards you, 
    at unprecedented speed.
    \item \textbf{The fallacy of mediocracy.} Iteration and failure are the method, not the excuse for settling.
    Do not confuse velocity with quality. Whatever you learn, whatever you build: make it good.
    \item \textbf{Outcomes emerge. They are not delivered.} Judgment, taste, the capacity to sit with uncertainty:
    these are not course outputs. They are what you build from the encounter.
    I can create the conditions. The rest depends on what you do with them.
  \end{itemize}
}

%%
% Page header: keep tufte-book's native two-side alternation
% (verso LE shows the book title, recto RO shows the current chapter
% title via \chaptermark/\leftmark) and add the author name to the
% verso only, so the author surfaces on every other page without
% disturbing the chapter title on recto.
%%
\let\tufte@orig@mainmatter\mainmatter
\renewcommand{\mainmatter}{%
  \tufte@orig@mainmatter
  \fancyhead[LE]{\thepage\quad\smallcaps{\newlinetospace{\plaintitle}}\quad\smallcaps{\defaultauthor}}%
}
%%
% Helper for the main file: call AFTER \mainmatter to add the author to
% the recto chapter header. Kept outside \renewcommand{\mainmatter} so
% that the inner \renewcommand{\chaptermark} sees its parameter as #1
% rather than ##1, which works cleanly with \markboth.
%%
\newcommand{\authorinrectoheader}{%
  \renewcommand{\chaptermark}[1]{\markboth{##1\quad\smallcaps{\defaultauthor}}{}}%
}

%%
% Font options template
%%
\newcommand{\loadoptionalfonts}{%
  % Font options (uncomment if available)
  %\IfFileExists{bergamo.sty}{\usepackage[osf]{bergamo}}{}% Bembo
  %\IfFileExists{chantill.sty}{\usepackage{chantill}}{}% Gill Sans
  %\usepackage{microtype}
}